\section{Solid Mechanics Department}
\label{sec:solid_mechanics}

\subsection{Laboratories}
There are two main labs located at Teknikringen 8D: the \textbf{Solid Mechanics Lab} and the \textbf{Lightweight Structures Lab}. To use these facilities, you must complete a 3-hour safety course and receive an introduction from the lab managers.

\subsection{KTH Service and Logistics}
The KTH Service team at Teknikringen 8D handles various administrative and logistical tasks, including parcel deliveries. If you have incoming parcels, please use the following address:

\begin{quote}
    KTH Verksamhetsstöd\\
    Skolan för teknikvetenskap\\
    Teknikringen 8D, 100 44 Stockholm\\
    ATT: [YOUR NAME]
\end{quote}

\subsection{Academic Requirements and Courses}
As a PhD student in this department, you are required to obtain \textbf{60 credits} to defend your thesis.
\begin{itemize}
    \item \textbf{Credit Transfer:} Up to 15 credits can typically be transferred from your Master's degree. This requires submitting a credit transfer form to the Doctoral Education Support (FOA).
    \item \textbf{Mandatory Courses:} Essential courses include Research Ethics, Continuum Mechanics, and Testing Techniques.
    \item \textbf{Planning:} Note that some departmental courses are only offered every 2 or 3 years, so plan your Individual Study Plan (ISP) carefully.
    \item \textbf{Course Levels:} Courses starting with `3` (e.g., 3000-series) are doctoral level, while `2` and `1` indicate Master's and Bachelor's levels respectively.
\end{itemize}

For the credit transfer form and more information, visit: 
\url{https://intra.kth.se/en/eecs/forskarutbildning/courses/course-information-for-doctoral-students-1.813447}

\subsection{Learning Swedish}
If you wish to learn Swedish, several options are available:
\begin{itemize}
    \item \textbf{SFI (Swedish for Immigrants):} Free courses provided by the municipality.
    \item \textbf{KTH Courses:} Language courses specifically for students.
    \item \textbf{Swedish for Employees:} You can request to join these through your supervisor.
\end{itemize}

\subsection{Lunch and Social Life}
Lunch is a central social event, typically occurring between 11:30 and 13:00. Most people gather in the second-floor lounge room (where the microwaves are) around 12:00.

\subsubsection{Dining Options Near Campus}
\begin{itemize}
    \item \textbf{Syster o Bror:} A popular nearby restaurant with a buffet.
    \item \textbf{Q Building:} Offers a large lunch buffet.
    \item \textbf{Brazilia:} High-quality buffet with many vegan options.
    \item \textbf{7-Eleven:} Convenient for quick snacks (use the \textit{Karma} app for discounts).
    \item \textbf{T-Snabben:} A minimarket with a salad bar (\textit{Picadeli}).
    \item \textbf{Coffice:} Good for coffee, fika, and sandwiches.
    \item \textbf{Nymble (Student Union Building):} Offers meals with student discounts.
\end{itemize}

\textbf{Pancake Thursday:} It is a Swedish tradition to serve pea soup and pancakes on Thursdays; you will find this menu at many of the locations listed above.

\subsection{Lounge Room and Kitchen Facilities}
The lounge room is fully equipped with microwaves, refrigerators, and dishwashers. Free coffee is often available during breaks. 
\textit{Note:} Tap water in Sweden is of very high quality and perfectly safe to drink; there is no need to buy bottled water.
